\documentclass{article} % For LaTeX2e
\usepackage{iclr2026_conference,times}

% Optional math commands from https://github.com/goodfeli/dlbook_notation.
\input{math_commands.tex}

\usepackage{hyperref}
\usepackage{url}
\usepackage{graphicx}
\usepackage{booktabs}
\usepackage{amsmath}
\usepackage{amssymb}
\usepackage{amsthm}
\usepackage{algorithm}
\usepackage{algorithmic}
\usepackage{multirow}
\usepackage{subcaption}


\newtheorem{theorem}{Theorem}
\newtheorem{lemma}[theorem]{Lemma}
\newtheorem{corollary}[theorem]{Corollary}
\newtheorem{proposition}[theorem]{Proposition}
\newtheorem{remark}{Remark}
\newtheorem{definition}{Definition}

\title{Probabilistic Decoupling for Nonlinear Inverse Problems: Overcoming Model Collapse in Unidentifiable Systems}

% Authors must not appear in the submitted version. They should be hidden
% as long as the \iclrfinalcopy macro remains commented out below.
\author{Anonymous Authors \\
Paper Under Review for ICLR 2026 \\
Submission ID: Click to Enter OpenReview ID
}

\newcommand{\fix}{\marginpar{FIX}}
\newcommand{\new}{\marginpar{NEW}}
\newcommand{\xobs}{\mX_{\mathrm{obs}}} % shorthand used in the theory sections

%\iclrfinalcopy % Uncomment for camera-ready version, but NOT for submission.
\begin{document}


\maketitle

\begin{abstract}
Estimating the parameters of nonlinear dynamical systems from sparse, partial, and noisy observations is a severely ill-posed inverse problem, and under structural non-identifiability it exposes a failure that is easy to mistake for mere inaccuracy: a network trained to minimize squared error collapses to the conditional mean $\E[\vtheta \mid \mX_{\mathrm{obs}}]$, which for a non-identifiable ridge is a low-density, often unphysical point rather than a solution. We take this collapse as a starting point, not a conclusion. Since the non-identifiability lives entirely in the unobserved states, the target posterior $p(\vtheta \mid \mX_{\mathrm{obs}})$ is a mixture of sharp, well-posed conditionals $p(\vtheta \mid \mX_{\mathrm{obs}}, \mX_{\mathrm{hid}})$ blurred by hidden-state uncertainty, and the natural sampler for it alternates between a hidden-state conditional and a parameter conditional---the same alternating operator that a deterministic decoupled estimator runs, but whose deterministic fixed point is exactly the collapsed conditional mean. We keep the operator and change what it converges to: replacing each deterministic map with a stochastic kernel turns the iteration into a sampler of the joint posterior. We give two guarantees in this probabilistic regime that parallel the contraction-and-alignment of deterministic work---injected noise makes the chain uniformly ergodic, so a unique invariant law exists with no contraction assumption, and training the two conditionals as denoisers steers that law toward the true posterior with a controlled bound. On sparsely observed physiological and synthetic systems, the framework recovers structured, multimodal posteriors along non-identifiable manifolds without collapse, and we characterize the resulting coverage--sharpness trade-off honestly. Full statements and proofs appear in a companion manuscript.
\end{abstract}

\section{Introduction}
\label{sec:intro}
Inferring the governing parameters of nonlinear dynamical systems from sparse, noisy, and partially observed trajectories is a cornerstone challenge across mathematical biology, ecology, and physiology. The task is fundamentally hard for three compounding reasons: parameter estimation is an ill-posed inverse problem; \emph{sparse} sampling maximizes its ill-conditioning; and \emph{partial} observation induces structural non-identifiability, in which many parameter configurations produce the same observable trajectory.

A natural response is amortized regression: train a network on simulated $(\mX_{\mathrm{obs}}, \vtheta)$ pairs and predict $\vtheta$ directly, in the spirit of simulation-based inference \citep{cranmer2020frontier}. Under structural non-identifiability, however, the squared-error solution is the conditional mean $\E[\vtheta \mid \mX_{\mathrm{obs}}]$, and for a non-identifiable ridge that mean lies in a low-density---often physically invalid---region between the branches of the solution set. The estimator is then not merely imprecise; it answers the wrong question, reporting a confident point where the data support an extended manifold. The same pathology afflicts \emph{deterministic} decoupled estimators, which split the map $\mX_{\mathrm{obs}} \mapsto \vtheta$ through the hidden state as $\mX_{\mathrm{obs}} \mapsto \mX_{\mathrm{hid}} \mapsto \vtheta$ and iterate two networks to a fixed point: enforcing a global contraction (e.g., via spectral normalization \citep{miyato2018spectral}) secures the fixed point but pins it to the same conditional mean, and the rigidity throttles representational capacity. We treat this collapse (regression-to-the-mean) as the \emph{start} of the problem, not its resolution: it tells us to recover the posterior, not to sharpen a point.

Why is $p(\vtheta \mid \mX_{\mathrm{obs}})$ so poorly behaved, and what fixes it? The non-identifiability is located entirely in the missing states. Conditioning on the \emph{full} trajectory renders $\vtheta$ identifiable, so the conditional $p(\vtheta \mid \mX_{\mathrm{obs}}, \mX_{\mathrm{hid}})$ is sharp and well-posed; the ridge appears only after marginalizing the hidden states,
\begin{equation}
\label{eq:mixture}
    p(\vtheta \mid \mX_{\mathrm{obs}}) = \int p(\vtheta \mid \mX_{\mathrm{obs}}, \mX_{\mathrm{hid}})\, p(\mX_{\mathrm{hid}} \mid \mX_{\mathrm{obs}})\, d\mX_{\mathrm{hid}} .
\end{equation}
This is the probabilistic rationale for decoupling: separate a well-posed parameter conditional from the hidden-state uncertainty that blurs it, rather than attacking the blurred marginal directly.

The natural sampler for this mixture alternates between the two conditionals---drawing $\mX_{\mathrm{hid}} \sim p(\mX_{\mathrm{hid}} \mid \mX_{\mathrm{obs}}, \vtheta)$ and then $\vtheta \sim p(\vtheta \mid \mX_{\mathrm{obs}}, \mX_{\mathrm{hid}})$---a (pseudo-)Gibbs sweep defining an operator $\mathcal{T}$. A deterministic decoupled estimator runs exactly this operator with the conditional \emph{means}, so its fixed point is $\E[\vtheta \mid \mX_{\mathrm{obs}}]$: the collapsed center is not an artifact of a particular architecture but the fixed point of the mean operator, a one-line consequence of iterating $\mathcal{T}$ with means. Our move is to keep the operator and change what it converges to. Replacing each deterministic map by a stochastic kernel lets the same iteration sample the full conditional instead of its mean, so the question shifts from ``does $\mathcal{T}$ have a fixed point'' to ``what law does $\mathcal{T}$ converge to, and is it the right one.'' Deterministic decoupled estimators are thus a baseline class and a limiting case, not a prerequisite. We instantiate the two conditionals as Conditional Variational Autoencoders (a Dual CVAE) and alternate between them with a stochastic Pseudo-Gibbs routine that samples $\mX_{\mathrm{hid}}$ and $\vtheta$ given $\mX_{\mathrm{obs}}$.

Two questions then replace the single fixed-point question, and we answer them in parallel to the contraction-and-alignment guarantees of deterministic work. \emph{Existence}: the injected noise gives $\mathcal{T}$ a strictly positive transition density on a compact state space, making it uniformly ergodic, so a unique invariant law $\pi^\ast$ exists and is reached geometrically---with \emph{no} contraction assumption, in contrast to the spectral-norm contraction a deterministic fixed point requires. \emph{Correctness}: training the two conditionals as denoisers steers $\pi^\ast$ toward the true posterior, the probabilistic counterpart of teacher forcing, with a bound that is controlled rather than zero and a self-consistency certificate computable from the trained model alone. We state these results in abbreviated form and defer full statements, assumptions, and proofs to a companion manuscript \citep{adcvae-theory}.

% ---------------------------------------------------------------------------
% Fig 1 (teaser) design spec --- draw as a two-panel conceptual diagram over a
% shared 2-D parameter plane (axes S_I, sigma), truth = star on a grey
% non-identifiability fiber (S_I * sigma = C, a crescent through the star).
%   Panel (a) "Collapse to the mean": single-net / deterministic decoupled.
%     Show a tight, overconfident cluster (near-point) parked at the conditional
%     mean E[theta|X_obs], visibly OFF the fiber, in the low-density interior
%     between fiber branches; a small caption tag "E[theta|X_obs]". Optional
%     faint arrow: iterate T (means) -> single fixed point.
%   Panel (b) "Track the fiber": Dual CVAE. Show a crescent-shaped point cloud
%     hugging the fiber and covering the star; a few Pseudo-Gibbs sample arrows
%     alternating along X_hid <-> theta to convey the sampler.
%   Shared legend: star = truth, grey curve = fiber, dots = posterior samples.
%   Keep it schematic; a real posterior inset from figure3.py (panel C) may be
%   dropped in later. Replace the \fbox below with the final PDF/TikZ.
% ---------------------------------------------------------------------------
\begin{figure}[h]
\begin{center}
\fbox{\rule[-.5cm]{0cm}{3.5cm} \rule[-.5cm]{5cm}{0cm}}
\end{center}
\caption{\textbf{Teaser: collapse to the mean vs.\ tracking the fiber.} Under partial observation the parameter posterior concentrates on a non-identifiable fiber (grey curve) passing through the truth (star). \textbf{(a)} A single-network or deterministic decoupled estimator, trained to minimize squared error, collapses to the conditional mean $\E[\vtheta \mid \mX_{\mathrm{obs}}]$---a confident point that sits off the fiber in a low-density region. \textbf{(b)} Our Dual CVAE runs the \emph{same} alternating operator with stochastic kernels, so the iteration samples the full conditional and traces the crescent-shaped fiber, covering the truth with calibrated, structured uncertainty.}
\label{fig:teaser}
\end{figure}

Our contributions are summarized as follows:
\begin{itemize}
    \item A \emph{generative decoupling} framework that preserves the physical hidden-state/parameter partition and recasts non-identifiability as structured posterior uncertainty; the two-network alternation is \emph{derived} as a sampler of the joint posterior via the mixture identity~\eqref{eq:mixture}, not adopted as a heuristic.
    \item Two guarantees that replace global spectral contraction with local, noise-driven ones: existence of a unique invariant law by ergodicity (assuming no contraction) and steering of that law toward the true posterior by denoising, accompanied by a computable self-consistency certificate. We state them briefly here and prove them in the companion manuscript \citep{adcvae-theory}.
    \item Empirical evidence, on a physiological model (the OGTT minimal model) and controlled synthetic systems, that the framework recovers structured multimodal posteriors along non-identifiable manifolds without collapse, while single-network and deterministic decoupled baselines collapse to the conditional mean; we characterize the coverage--sharpness trade-off honestly rather than claiming uniform superiority.
\end{itemize}

\section{Background \& Related Work}
\label{sec:background}
Consider a nonlinear dynamical system governed by the ordinary differential equation (ODE) $\dot{\vx}(t) = f(\vx(t); \vtheta)$, where $\vx(t) \in \mathbb{R}^d$ represents the full system states and $\vtheta \in \mathbb{R}^p$ denotes the constant parameters. In realistic scenarios, we only observe a subset of the states $\vx_{\mathrm{obs}} \subset \vx$ at highly sparse time intervals, leaving $\vx_{\mathrm{hid}} = \vx \setminus \vx_{\mathrm{obs}}$ latent.

A \emph{deterministic decoupled} baseline instantiates the two maps as networks $\mX_{\mathrm{hid}}^{(k+1)} = H_\phi(\mX_{\mathrm{obs}}, \vtheta^{(k)})$ and $\vtheta^{(k+1)} = P_\psi(\mX_{\mathrm{obs}}, \mX_{\mathrm{hid}}^{(k+1)})$, defining the sweep operator $\mathcal{T}(\vtheta; \mX_{\mathrm{obs}}) = P_\psi(\mX_{\mathrm{obs}}, H_\phi(\mX_{\mathrm{obs}}, \vtheta))$. Enforcing a global Lipschitz constant $\mathrm{Lip}(\mathcal{T}) < 1$ via spectral normalization \citep{miyato2018spectral} makes $\mathcal{T}$ a contraction with a unique Banach fixed point. Because the two networks are trained to output conditional means, that fixed point is the conditional mean $\E[\vtheta \mid \mX_{\mathrm{obs}}]$: under structural non-identifiability it lands in the low-density interior of the solution ridge, and the contraction constraint additionally bottlenecks capacity. Our framework retains the operator $\mathcal{T}$ but replaces each deterministic map with a stochastic transition kernel (\Secref{sec:method}), so the same iteration samples the joint posterior instead of contracting to its mean.

\subsection{Related Work}
\label{sec:related}

\paragraph{Simulation-based inference.}
Amortized posterior estimation trains a conditional density on simulated parameter--data pairs and reuses it across observations \citep{cranmer2020frontier}. Neural posterior estimation and its sequential refinements \citep{papamakarios2016fast, greenberg2019automatic}, benchmarked systematically by \citet{lueckmann2021benchmarking}, typically fit a single conditional $p(\vtheta \mid \mX_{\mathrm{obs}})$ with a normalizing flow \citep{papamakarios2021normalizing}. We differ by \emph{decoupling} through the physical hidden state: rather than model the blurred marginal directly, we sample the sharp conditional $p(\vtheta \mid \mX_{\mathrm{obs}}, \mX_{\mathrm{hid}})$ and marginalize $\mX_{\mathrm{hid}}$ by iteration, injecting ODE structure that a single flow must otherwise learn implicitly.

\paragraph{Score-based generative models.}
That target-noise training recovers the score of the data distribution \citep{vincent2011connection, efron2011tweedie} underpins modern score-based and diffusion models \citep{song2019generative, ho2020denoising, song2021score}, as well as flow-matching and stochastic-interpolant alternatives \citep{lipman2023flow}. Our steering guarantee uses exactly this denoiser-as-score correspondence, but applied to the two learned \emph{conditionals} of the inverse problem rather than to a single unconditional prior.

\paragraph{Score-based posterior sampling and split Gibbs.}
The closest modern neighbors sample a posterior by alternating a data-consistency step with a learned score/denoiser prior: diffusion posterior sampling \citep{chung2023diffusion}, denoiser-driven stochastic solvers \citep{kadkhodaie2021stochastic}, and split-Gibbs / plug-and-play samplers \citep{vono2019split, coeurdoux2024plug}. These target a \emph{known} likelihood paired with a generic prior; our two blocks are instead \emph{two learned conditionals of a partially observed dynamical system}, coupled through the hidden state, and we quantify the bias from their possible incompatibility with a computable certificate.

\paragraph{Compatible conditionals and dependency networks.}
Sampling from independently learned conditionals is pseudo-Gibbs on a dependency network \citep{heckerman2000dependency}, whose stationary law need not correspond to any genuine joint unless the conditionals are compatible \citep{arnold2001conditionally}. Prior treatments are discrete; we extend the construction to a continuous, non-identifiable state space and, instead of assuming compatibility, measure the residual incompatibility and bound its effect (\Secref{subsec:theory}).

\section{Problem Formulation}
\label{sec:problem}

\subsection{Dynamical System Modeling via ODEs}
Dynamical system modeling provides a mathematical framework for analyzing time-varying phenomena
across the physical, engineering, and biological sciences. Such systems are typically governed by
a system of ordinary differential equations (ODEs),
\begin{equation}
\label{eqn:ode}
    \dot{\vx}(t) = f\big(\vx(t); \vtheta\big), \qquad
    \vx(t) = \begin{bmatrix} \vx_{\mathrm{obs}}(t)\\ \vx_{\mathrm{hid}}(t)\end{bmatrix} \in \mathbb{R}^{d},
\end{equation}
where $\vx(t)$ is the state vector at time $t$, $\vtheta \in \mathbb{R}^{p}$ denotes the
time-invariant parameters governing the dynamics, and $f$ is a sufficiently smooth vector field
ensuring existence and uniqueness of the solution trajectory.

In classical system identification, one estimates the true parameters $\vtheta$ by observing the
state trajectory $\vx(t)$ over an interval $[0,T]$. Once $\vtheta$ is identified, the physical
properties of the system can be analyzed and simulated. In practice, however, full and continuous
measurement of $\vx(t)$ is rarely available. We formalize this through two distinct bottlenecks.

\paragraph{Partial Observability (Non-identifiability).}
The full state is rarely observable in its entirety; it is partitioned as
\begin{equation}
    \vx(t) = \begin{bmatrix} \vx_{\mathrm{obs}}(t) \\ \vx_{\mathrm{hid}}(t) \end{bmatrix},
    \quad \vx_{\mathrm{obs}}(t) \in \mathbb{R}^{d_{\mathrm{obs}}},
    \quad \vx_{\mathrm{hid}}(t) \in \mathbb{R}^{d_{\mathrm{hid}}},
    \quad d = d_{\mathrm{obs}} + d_{\mathrm{hid}},
\end{equation}
where $\vx_{\mathrm{obs}}$ collects the observable states and $\vx_{\mathrm{hid}}$ the unobserved
(hidden) ones. For instance, in the Oral Glucose Tolerance Test (OGTT) a physiological model
couples blood glucose and insulin ($\vx = [x_{\mathrm{glucose}}, x_{\mathrm{insulin}}]^\top$), yet
clinical protocols often measure glucose only ($\vx_{\mathrm{obs}} = x_{\mathrm{glucose}}$). The
absence of $\vx_{\mathrm{hid}}$ induces severe \textit{non-identifiability}: infinitely many
configurations of $\vtheta$ can yield an identical observable trajectory, preventing deterministic
optimizers from converging to the true parameter.

\paragraph{Sparse Sampling (Temporal Aliasing).}
Continuous monitoring is likewise impractical; states are collected only at discrete, sparse times
$\mathcal{T} = \{t_0, t_1, \dots, t_N\}$ with $0 = t_0 < \dots < t_N = T$, yielding
\begin{equation}
    \mX_{\mathrm{obs}} = \big[\, \vx_{\mathrm{obs}}(t_0),\, \dots,\, \vx_{\mathrm{obs}}(t_N) \,\big]
    \in \mathbb{R}^{d_{\mathrm{obs}} \times (N+1)}.
\end{equation}
Sparse sampling introduces temporal aliasing, making rapidly oscillating trajectories
indistinguishable. Moreover, classical pipelines that rely on numerical differentiation
($\dot{\vx}(t) \approx [\vx(t+\Delta t) - \vx(t)]/\Delta t$) incur devastating approximation errors,
rendering gradient-based dynamics learning highly unstable.

\subsection{Goal: Joint Posterior Estimation via Coupled Generative Models}
To estimate $\vtheta$ under partial and sparse observation, we recast parameter estimation as
probabilistic inference: our goal is to characterize the posterior
$p(\vtheta \mid \mX_{\mathrm{obs}})$ given only the sparse, partial data $\mX_{\mathrm{obs}}$.

Because $\mX_{\mathrm{hid}}$ is entirely absent, directly optimizing
$p(\vtheta \mid \mX_{\mathrm{obs}})$ is ill-posed and numerically intractable. We therefore model
the \textit{joint posterior} over both the parameters and the missing hidden trajectory,
$p(\vtheta, \mX_{\mathrm{hid}} \mid \mX_{\mathrm{obs}})$, where
\begin{equation}
    \mX_{\mathrm{hid}} = \big[\, \vx_{\mathrm{hid}}(t_0),\, \dots,\, \vx_{\mathrm{hid}}(t_N) \,\big]
    \in \mathbb{R}^{d_{\mathrm{hid}} \times (N+1)}.
\end{equation}
We decompose this joint estimation into two coupled conditional generative models, parameterized by
$\phi$ and $\psi$:
\begin{align}
    \text{\textbf{Hidden Trajectory Generator:}} \quad & p_\phi(\mX_{\mathrm{hid}} \mid \mX_{\mathrm{obs}}, \vtheta), \\
    \text{\textbf{Parameter Generator:}} \quad & p_\psi(\vtheta \mid \mX_{\mathrm{obs}}, \mX_{\mathrm{hid}}).
\end{align}
Given an initial estimate $\vtheta^{(0)}$, we approximate the intractable joint posterior by an
alternating sampling scheme, for $k = 0, 1, 2, \dots$:
\begin{align}
    \mX_{\mathrm{hid}}^{(k+1)} & \sim p_\phi(\mX_{\mathrm{hid}} \mid \mX_{\mathrm{obs}}, \vtheta^{(k)}), \label{eq:update_x} \\
    \vtheta^{(k+1)} & \sim p_\psi(\vtheta \mid \mX_{\mathrm{obs}}, \mX_{\mathrm{hid}}^{(k+1)}). \label{eq:update_theta}
\end{align}
This process refines the uncertainty in both the hidden-state and parameter spaces simultaneously,
guiding the system toward feasible parameters rather than a single point estimate.

\paragraph{Amortized Inference via Synthetic Simulation.}
To avoid per-instance MCMC or online grid search, we adopt an \textit{amortized} framework. We
assume access to a synthetic dataset
$\mathcal{D} = \{ (\mX_{\mathrm{obs}}^{(m)}, \mX_{\mathrm{hid}}^{(m)}, \vtheta^{(m)}) \}_{m=1}^{M}$,
generated by sampling parameters $\vtheta^{(m)}$ and initial conditions $\vx^{(m)}(0)$ from a
physically plausible prior $p(\vtheta, \vx_0)$ and integrating the forward ODE,
\begin{equation}
    \vx^{(m)}(t) = \vx^{(m)}(0) + \int_{0}^{t} f\big(\vx^{(m)}(\tau); \vtheta^{(m)}\big)\, d\tau.
\end{equation}
Training $p_\phi$ and $p_\psi$ on $\mathcal{D}$ amortizes the inference cost into an offline phase;
once trained, the networks approximate the joint posterior and enable instantaneous,
distribution-aware recovery for any newly observed sparse sequence $\mX_{\mathrm{obs}}$.

\section{Proposed Method: Alternating Denoising CVAEs (A-DCVAE)}
\label{sec:method}

To address structural non-identifiability in sparsely observed dynamical systems, we decompose the
intractable joint posterior $p(\vtheta, \mX_{\mathrm{hid}} \mid \mX_{\mathrm{obs}})$ into two coupled
conditional sub-tasks, modeled by two Conditional Variational Autoencoders \citep[CVAEs;][]{kingma2014auto, sohn2015learning}: the
\textit{Hidden-State CVAE} $\gH_\phi$ and the \textit{Parameter CVAE} $\gP_\psi$. Inference
alternates between them in the spirit of Gibbs sampling. We refer to the resulting framework as
\textbf{Alternating Denoising CVAEs (A-DCVAE)}.

\subsection{Architecture and Denoising Objective}
\label{subsec:architecture}

Each CVAE introduces a continuous latent variable, $\vz_H$ and $\vz_P$, with a standard isotropic
Gaussian prior $\mathcal{N}(\vzero, \mI)$. The generative decoders $g_\phi$ and $g_\psi$ are
deterministic maps that output the \textit{mean} of the target variable. During training, the
amortized encoders $q_\phi$ and $q_\psi$ map ground-truth triplets
$(\mX_{\mathrm{obs}}, \mX_{\mathrm{hid}}, \vtheta)$ to their latent spaces, and each decoder is
trained to reconstruct its clean target.

\paragraph{Dual noise injection.}
To robustify the inverse mappings against regression-to-the-mean, we cast both CVAEs as
\textit{conditional denoising autoencoders}. We inject independent Gaussian noise into the
\textit{targets} ($\boldsymbol{\epsilon}_t \sim \mathcal{N}(\vzero, \sigma_t^2 \mI)$) and the
\textit{conditions} ($\boldsymbol{\epsilon}_c \sim \mathcal{N}(\vzero, \sigma_c^2 \mI)$), and define
the perturbed variables
\begin{equation}
    \tilde{\mX}_{\mathrm{hid}} = \mX_{\mathrm{hid}} + \boldsymbol{\epsilon}_t, \quad
    \tilde{\vtheta} = \vtheta + \boldsymbol{\epsilon}_t, \quad
    \check{\mX}_{\mathrm{hid}} = \mX_{\mathrm{hid}} + \boldsymbol{\epsilon}_c, \quad
    \check{\vtheta} = \vtheta + \boldsymbol{\epsilon}_c .
\end{equation}
The encoder receives the \emph{target}-perturbed variable while the decoder is conditioned on the
\emph{condition}-perturbed variable, and both are penalized to reconstruct the \emph{clean} target.
This yields the $\beta$-annealed denoising ELBO objectives (minimized), which decouple across the
two networks:
\begin{align}
\label{eqn:loss_H}
\gL_H(\phi) &= \E_{\boldsymbol{\epsilon}_t, \boldsymbol{\epsilon}_c}
    \E_{q_\phi(\vz_H \mid \mX_{\mathrm{obs}}, \tilde{\mX}_{\mathrm{hid}}, \check{\vtheta})}
    \Big[ \big\| \mX_{\mathrm{hid}} - g_\phi(\mX_{\mathrm{obs}}, \check{\vtheta}, \vz_H) \big\|_F^2 \Big]
    + \beta\, D_{\mathrm{KL}}\!\big( q_\phi \,\|\, p(\vz_H) \big), \\
\label{eqn:loss_P}
\gL_P(\psi) &= \E_{\boldsymbol{\epsilon}_t, \boldsymbol{\epsilon}_c}
    \E_{q_\psi(\vz_P \mid \mX_{\mathrm{obs}}, \check{\mX}_{\mathrm{hid}}, \tilde{\vtheta})}
    \Big[ \big\| \vtheta - g_\psi(\mX_{\mathrm{obs}}, \check{\mX}_{\mathrm{hid}}, \vz_P) \big\|_2^2 \Big]
    + \beta\, D_{\mathrm{KL}}\!\big( q_\psi \,\|\, p(\vz_P) \big),
\end{align}
where $\beta$ is annealed from $0$ to $\beta_{\max}$ over a warmup phase to prevent posterior
collapse. As we show in \Secref{subsec:theory}, the two noise channels play distinct and
complementary roles: \emph{condition noise} controls stability, while \emph{target noise} controls
directional correctness.

\subsection{Test-Time Inference via Stochastic Pseudo-Gibbs Sampling}
\label{subsec:inference}

At test time, sampling from the joint posterior is intractable. \citet{heckerman2000dependency}
showed that \textit{pseudo-Gibbs sampling} over independently learned---and possibly
inconsistent---conditional distributions still defines an ergodic chain converging to a stationary
joint distribution, but their analysis is confined to discrete domains. Naively applying pseudo-Gibbs
to \emph{deterministic} neural decoders in a continuous, non-identifiable domain collapses the chain
to a single, physically meaningless mean (\Secref{subsec:theory}, Issue~3).

\paragraph{Stochastic transition operators via infinite Gaussian mixtures.}
We instead construct stochastic transition kernels by injecting scaled Gaussian noise into the
deterministic decoder outputs and marginalizing over the latent prior. At iteration $k$,
\begin{align}
p(\mX_{\mathrm{hid}}' \mid \mX_{\mathrm{obs}}, \vtheta)
    &= \int \mathcal{N}\!\big(\mX_{\mathrm{hid}}' ;\, g_\phi(\mX_{\mathrm{obs}}, \vtheta, \vz_H),\, \sigma_H^2 \mI\big)\, \mathcal{N}(\vz_H ;\, \vzero, \mI)\, d\vz_H, \\
p(\vtheta' \mid \mX_{\mathrm{obs}}, \mX_{\mathrm{hid}}')
    &= \int \mathcal{N}\!\big(\vtheta' ;\, g_\psi(\mX_{\mathrm{obs}}, \mX_{\mathrm{hid}}', \vz_P),\, \sigma_P^2 \mI\big)\, \mathcal{N}(\vz_P ;\, \vzero, \mI)\, d\vz_P.
\end{align}
Drawing a fresh $\vz$ and adding Gaussian noise at every step realizes an \textit{infinite mixture of
Gaussians}, endowing the continuous Markov chain with the capacity to represent the highly complex,
multi-modal probability valleys of the unidentifiable ODE space. The complete transition operator
from $(\mX_{\mathrm{hid}}, \vtheta)$ to $(\mX_{\mathrm{hid}}', \vtheta')$ is
\begin{equation}
\label{eqn:kernel}
    \mathcal{T}\big((\mX_{\mathrm{hid}}, \vtheta) \to (\mX_{\mathrm{hid}}', \vtheta')\big)
    = p(\mX_{\mathrm{hid}}' \mid \mX_{\mathrm{obs}}, \vtheta)\; p(\vtheta' \mid \mX_{\mathrm{obs}}, \mX_{\mathrm{hid}}').
\end{equation}
The procedure is summarized in \Algref{alg:pseudo_gibbs}.

\begin{algorithm}[htbp]
\caption{Stochastic Pseudo-Gibbs Inference}
\label{alg:pseudo_gibbs}
\begin{algorithmic}[1]
\REQUIRE Sparse observation $\mX_{\mathrm{obs}}$; trained decoders $g_\phi, g_\psi$; noise scales $\sigma_H, \sigma_P$; steps $K$; burn-in $B$
\STATE Initialize $\vtheta^{(0)}$ (replicated across parallel chains)
\FOR{$k = 0$ \TO $K-1$}
\STATE Sample $\vz_H \sim \mathcal{N}(\vzero, \mI)$ and $\boldsymbol{\epsilon}_H \sim \mathcal{N}(\vzero, \sigma_H^2 \mI)$
\STATE $\mX_{\mathrm{hid}}^{(k+1)} \leftarrow g_\phi(\mX_{\mathrm{obs}}, \vtheta^{(k)}, \vz_H) + \boldsymbol{\epsilon}_H$
\STATE Sample $\vz_P \sim \mathcal{N}(\vzero, \mI)$ and $\boldsymbol{\epsilon}_P \sim \mathcal{N}(\vzero, \sigma_P^2 \mI)$
\STATE $\vtheta^{(k+1)} \leftarrow g_\psi(\mX_{\mathrm{obs}}, \mX_{\mathrm{hid}}^{(k+1)}, \vz_P) + \boldsymbol{\epsilon}_P$
\STATE (Optional) apply physical boundary projections $\mathcal{B}_y, \mathcal{B}_p$
\ENDFOR
\RETURN samples $\{(\mX_{\mathrm{hid}}^{(k)}, \vtheta^{(k)})\}_{k=B}^{K}$ after discarding burn-in $B$
\end{algorithmic}
\end{algorithm}

\subsection{Theoretical Properties: Three Issues, Three Guarantees}
\label{subsec:theory}

Running pseudo-Gibbs with deep networks in a continuous, non-identifiable space faces three
failure modes, each resolved by one component of A-DCVAE (proofs in \Secref{app:proofs}).

\paragraph{Issue 1 (instability / error amplification).}
Even when trained on exact data, the two CVAEs are only \textit{$\epsilon$-incompatible}: limited
network capacity and variational error mean no single joint distribution induces both conditionals
exactly. Because the manifold of physically valid trajectories $\mathcal{M}$ is a thin,
measure-zero set, an unguarded transition can amplify a single small step error exponentially through
the feedback loop, causing divergence.

\begin{theorem}[Condition noise induces local Jacobian regularization]
\label{thm:contraction}
Injecting condition noise $\boldsymbol{\epsilon}_c$ and reconstructing the clean target is, as
$\sigma_c \to 0$, asymptotically equivalent to a local Tikhonov penalty on the Frobenius norm of the
decoder's condition-Jacobian:
\begin{equation}
    \E_{\boldsymbol{\epsilon}_c}\big[\, \| g(\, \cdot + \boldsymbol{\epsilon}_c) - Y \|^2 \,\big]
    \;\approx\; \| g(\,\cdot\,) - Y \|^2 \;+\; \sigma_c^2\, \big\| \nabla g(\,\cdot\,) \big\|_F^2 .
\end{equation}
Suppressing $\|\nabla g\|_F$ smooths the conditional map and imposes strong dissipativity (a
Lyapunov drift) on the transition operator, so the chain remains bounded and stable rather than
diverging.
\end{theorem}

\paragraph{Issue 2 (wrong direction).}
Stability alone does not guarantee correctness: a stable chain may still settle at a spurious mean.
The chain must be attracted toward the true data manifold.

\begin{theorem}[Target noise induces denoising score matching]
\label{thm:score}
Following the denoising--score-matching equivalence \citep{vincent2011connection}, the global optimum of the target-denoising objective satisfies Tweedie's identity,
\begin{equation}
    g^*(\tilde{Y}) \;\approx\; \tilde{Y} + \sigma_t^2\, \nabla_Y \log p\big(Y \mid \mX_{\mathrm{obs}}, \vtheta\big),
\end{equation}
so the decoder's residual equals the (smoothed) score of the conditional data manifold. The update
therefore performs a Langevin drift toward high-density, physically consistent regions, forming an
attractor field that pulls off-manifold samples back onto $\mathcal{M}$.
\end{theorem}

\paragraph{Issue 3 (mode collapse).}
Passing only decoder means between the networks makes regression-to-the-mean inevitable in
non-identifiable regimes. The stochastic, infinite-mixture kernel~\eqref{eqn:kernel} removes this
failure.

\begin{theorem}[Ergodicity and score-aligned convergence]
\label{thm:ergodicity}
Under the dissipativity of Theorem~\ref{thm:contraction} and compact boundary projections, the infinite
Gaussian mixture kernel $\mathcal{T}$ is geometrically ergodic in the $2$-Wasserstein metric and
admits a unique invariant distribution $\pi^*$. Because $\mathcal{T}$ has the capacity to represent
multi-modal (e.g.\ crescent-shaped) supports, and the Langevin drift of Theorem~\ref{thm:score} provides
a persistent attraction toward the data manifold, the chain bypasses mode collapse and $\pi^*$ aligns
with the true joint posterior $p_{\mathrm{data}}(\vtheta, \mX_{\mathrm{hid}} \mid \mX_{\mathrm{obs}})$.
\end{theorem}

Together, Theorems~\ref{thm:contraction}--\ref{thm:ergodicity} replace the rigid \emph{global} spectral
constraint of prior deterministic work with a softer, data-driven \emph{local} regularization that
preserves expressiveness while guaranteeing a stable, correctly directed, collapse-free chain.




\section{Experimental Results: Identifiability \& Model Collapse}
\label{sec:experiments}

We evaluate our framework on physiological models known to contain structurally unidentifiable regimes under sparse conditions, notably the Oral Glucose Tolerance Test (OGTT) minimal model. Trajectories are generated within a clinically valid compact support representing metabolic variants from normal to diabetic ranges.

\begin{figure}[h]
\begin{center}
\fbox{\rule[-.5cm]{0cm}{3.5cm} \rule[-.5cm]{3.5cm}{0cm}} \quad
\fbox{\rule[-.5cm]{0cm}{3.5cm} \rule[-.5cm]{3.5cm}{0cm}}
\end{center}
\caption{\textbf{MCMC Convergence and Posterior Quality:} (Left) Trace plots demonstrating stable empirical convergence and mixing after a brief burn-in period. (Right) Autocorrelation metrics confirming rapid decay under local Jacobian constraints.}
\label{fig:mcmc_diag}
\end{figure}

\subsection{Overcoming Overconfidence and Capturing Physical Non-Identifiability}
\label{sec:non_identifiability}

A fundamental challenge in inverse problems for physiological systems is the inherent non-identifiability of parameters; multiple parameter configurations can yield identical observational data. An ideal probabilistic model must not only provide accurate point estimates but also reliably quantify this epistemic uncertainty by mapping the true solution manifold.

As shown in Figure \ref{fig:posterior_analysis}a, both the Single CVAE baseline and the proposed A-DCVAE appear to capture the global population distribution $p(\vtheta)$ reasonably well. This global coverage explains the competitive point-estimate metrics of the baseline model (Table \ref{tab:quantitative_results}), where the Single CVAE achieves nominally lower RMSE for both $S_I$ and $\sigma$. However, relying solely on standard regression metrics like RMSE in non-identifiable systems yields an ``illusion of accuracy.'' 

\begin{figure}[t]
    \centering
    \begin{subfigure}[b]{0.48\textwidth}
        \centering
        \includegraphics[width=\textwidth]{figures/figure3_B_global_2D.pdf}
        \caption{Global Joint $p(\vtheta)$}
        \label{fig:global_joint}
    \end{subfigure}
    \hfill
    \begin{subfigure}[b]{0.48\textwidth}
        \centering
        \includegraphics[width=\textwidth]{figures/figure3_C_local_2D.pdf}
        \caption{Local Posterior $q(\vtheta \mid \mX_{\mathrm{obs}})$}
        \label{fig:local_posterior}
    \end{subfigure}
    \caption{Visualizing the probabilistic capabilities of the models. (\textbf{a}) At the population level, both models adequately cover the global prior distribution. (\textbf{b}) At the individual sample level, the Single CVAE (blue) exhibits severe pathological mode collapse, predicting an extremely narrow distribution that misses the ground truth (red star). In contrast, A-DCVAE (green) accurately model the theoretical non-identifiability manifold ($S_I \times \sigma = C$, purple dashed line), providing well-calibrated posterior uncertainty.}
    % figure3_B/C regenerated with unified "A-DCVAE" legend via
    % prob_models/paper/experiments/regen_figure3.py (canonical checkpoints, seed 42).
    \label{fig:posterior_analysis}
\end{figure}

The critical failure of the baseline becomes evident when examining the local posterior $q(\vtheta \mid \mX_{\mathrm{obs}})$ for a specific observation. As illustrated in the inset of Figure \ref{fig:posterior_analysis}b, the Single CVAE suffers from pathological mode collapse. It exhibits extreme overconfidence, predicting a nearly point-like distribution (MPIW of 0.0102) that completely misses the true parameter. This phenomenon indicates that the baseline network merely memorizes a regression-to-the-mean mapping rather than understanding the underlying data generation mechanism. Consequently, this overconfidence results in a catastrophically poor Prediction Interval Coverage Probability (PICP) of 17.20\% and an exploding Negative Log-Likelihood (NLL) of 31990.99, rendering the model's probabilistic outputs practically useless for clinical or scientific inference.

\begin{table}[b]
    \centering
    \caption{Quantitative evaluation of parameter estimation and uncertainty calibration.}
    \label{tab:quantitative_results}
    \resizebox{\textwidth}{!}{
    \begin{tabular}{lcccccc}
        \toprule
        \multirow{2}{*}{\textbf{Model}} & \multicolumn{2}{c}{\textbf{Point Estimation}} & \multicolumn{4}{c}{\textbf{Probabilistic Calibration}} \\
        \cmidrule(lr){2-3} \cmidrule(lr){4-7}
        & \textbf{RMSE ($S_I \downarrow$ / $\sigma \downarrow$)} & \textbf{Pearson r ($S_I \uparrow$ / $\sigma \uparrow$)} & \textbf{PICP ($\uparrow$)} & \textbf{MPIW ($\downarrow$)} & \textbf{CRPS ($\downarrow$)} & \textbf{NLL ($\downarrow$)} \\
        \midrule
        Single CVAE & \textbf{0.1452} / \textbf{0.1480} & \textbf{0.8957} / \textbf{0.9371} & 17.20\% & \textbf{0.0102} & \textbf{0.0561} & 31990.99 \\
        A-DCVAE (Ours) & 0.1642 / 0.1811 & 0.8879 / 0.9311 & \textbf{99.44\%} & 1.1466 & 0.0836 & \textbf{-0.8076} \\
        \bottomrule
    \end{tabular}
    }
\end{table}

In stark contrast, the proposed A-DCVAE demonstrates a profound capacity to learn the physical laws governing the system. As shown in Figure \ref{fig:posterior_analysis}b, the generated posterior seamlessly stretches along the theoretical non-identifiability curve ($S_I \times \sigma = C$), successfully encompassing the ground truth within its predicted high-density region. By fully acknowledging the inherent uncertainty of the inverse problem, our approach achieves near-perfect calibration with a 99.44\% PICP and a highly stable negative NLL (-0.8076). These results demonstrate that while deterministic-like point estimation may yield marginally better RMSE, A-DCVAE provides a mathematically sound and physically valid probabilistic framework essential for reliable parameter inference.

\paragraph{Coverage versus sharpness, read honestly.}
The 99.44\% PICP is bought with a far wider mean interval (MPIW $0.0102 \to 1.1466$), which invites the objection that our posterior is merely under-confident. Two points sharpen the reading. First, the comparison is not width-for-width: the Single CVAE's narrow interval is \emph{false} confidence---it misses the truth (PICP 17.20\%), so its small MPIW measures precision about the wrong location. Second, and more importantly, under genuine non-identifiability a wide marginal is the \emph{correct} answer only if the width is spent along the unidentifiable direction and not across the identifiable one; aggregate MPIW cannot make this distinction, since it collapses a two-dimensional interval into a single scalar. The appropriate test is directional---whether the posterior stays tight across the identifiable combination ($mDI = S_I \times \sigma$) while extending along the non-identifiable fiber---which we assess directly against an analytic reference posterior next (\Secref{sec:reference}) and against strong baselines in \Secref{sec:baselines}. Judged by interval \emph{shape} relative to the fiber rather than by scalar width, A-DCVAE is faithful exactly where the single baseline is falsely confident.

\subsection{Validation against an analytic reference posterior}
\label{sec:reference}

Because glucose depends on $\vtheta$ only through the product $mDI = S_I\,\sigma$, the OGTT posterior admits a cheap gold standard: $p(S_I,\sigma \mid G) \propto \mathcal{L}(mDI)\,p(S_I,\sigma)$, evaluated on a grid by forward-simulation. This lets us grade the sampler against a reference rather than by self-consistency alone. Figure~\ref{fig:reference} overlays A-DCVAE's stationary samples $\pi^\ast$ on the reference density. Aggregated over eight observations, the reference confirms the identifiability structure numerically: its MAP $mDI$ is within $2.4\% \pm 2.0\%$ of the truth and its across-fiber coefficient of variation in $mDI$ is $10.0\% \pm 1.4\%$ (tight $\Rightarrow$ $mDI$ is identifiable), while the along-fiber extent is set by the prior. A-DCVAE follows the fiber \emph{direction} but is over-dispersed: its across-fiber $mDI$ CV is $80.5\% \pm 3.5\%$ against the reference's $10\%$, with a consistent offset (sliced-$W_2$ $0.20 \pm 0.14$; $W_1$ in $\log mDI$ $0.489 \pm 0.011$). We report this honestly as the visible approximation gap---the $\varepsilon_{\mathrm{app}}$/$\varepsilon_{\mathrm{inc}}$ of \Secref{subsec:theory} made concrete: the method recovers the correct one-dimensional support but not yet its exact width, an over-coverage rather than the baseline's catastrophic miss.

\begin{figure}[t]
    \centering
    \includegraphics[width=0.62\textwidth]{figures/figure_b3_reference_posterior.pdf}
    \caption{\textbf{Analytic reference posterior vs.\ A-DCVAE.} The reference density $p(S_I,\sigma \mid G) \propto \mathcal{L}(mDI)\,p(S_I,\sigma)$ (color) is a thin crescent on the fiber $S_I\,\sigma = mDI$ (dashed). A-DCVAE's samples $\pi^\ast$ (green) follow the fiber direction through the ground truth (star) but are over-dispersed across it---the honest approximation gap.}
    \label{fig:reference}
\end{figure}

\subsection{Does decoupling help? A matched-budget baseline comparison}
\label{sec:baselines}

A natural worry is that decoupling into two conditionals buys nothing over directly modeling $p(\vtheta \mid \mX_{\mathrm{obs}})$ with a single expressive network. We test A-DCVAE against a self-contained conditional-flow NPE baseline (a single RealNVP over $p(\vtheta \mid \mX_{\mathrm{obs}})$) and a deterministic MSE regressor (a point estimator), all trained on an identical context (same data, split, and normalizer) per run and graded against the analytic reference ($5$ seeds $\times$ $5$ observations). Table~\ref{tab:baselines} and Figure~\ref{fig:baselines} report the canonical-scale ($50$k-simulation, patience-matched) result. A-DCVAE is closest to the reference in sliced-$W_2$ ($0.170$ vs.\ the flow's $0.533$) and competitive in direction ($mDI$ rel-err $0.016$ vs.\ $0.013$), while the deterministic regressor collapses (along-fiber std exactly $0$). The flow's much higher raw HPD-$95$ coverage ($0.57$ vs.\ $0.20$) is not accuracy but over-dispersion: its along-fiber spread ($2.10$) is roughly three times both the reference scale and A-DCVAE's ($0.71$).

\begin{table}[t]
    \centering
    \caption{Canonical-scale ($50$k, patience-matched) comparison against the analytic reference posterior, mean $\pm$ std over $5$ seeds $\times$ $5$ observations. Best in each column in bold.}
    \label{tab:baselines}
    \begin{tabular}{lcccc}
        \toprule
        \textbf{Method} & \textbf{ref-HPD95 cov.} & \textbf{along-fiber std} & \textbf{sliced-$W_2$ ($\downarrow$)} & \textbf{$mDI$ rel-err ($\downarrow$)} \\
        \midrule
        A-DCVAE (dual, ours) & $0.203 \pm 0.02$ & $0.712 \pm 0.01$ & $\mathbf{0.170 \pm 0.10}$ & $0.016 \pm 0.01$ \\
        NPE-flow (single RealNVP) & $0.572 \pm 0.06$ & $2.105 \pm 0.23$ & $0.533 \pm 0.07$ & $\mathbf{0.013 \pm 0.01}$ \\
        det-reg (point) & $1.000^{\dagger}$ & $0.000$ & $0.292 \pm 0.03$ & $0.020 \pm 0.01$ \\
        \bottomrule
    \end{tabular}
    \\[2pt]
    {\footnotesize $^{\dagger}$ A point estimate has a degenerate coverage of $1$ (or $0$); it carries no interval and is reported only for context.}
\end{table}

\begin{figure}[t]
    \centering
    \includegraphics[width=\textwidth]{figures/figure_b4_baselines.pdf}
    \caption{\textbf{Does decoupling help?} (Left) posterior samples for one observation against the reference fiber; (right) the four metrics of Table~\ref{tab:baselines}. A-DCVAE (teal) is closest to the reference in sliced-$W_2$; NPE-flow (blue) is over-dispersed (high raw coverage, wide along-fiber spread); det-reg (red) collapses to a point.}
    \label{fig:baselines}
\end{figure}

We flag one honest caveat: at a smaller training scale ($10$k) the sliced-$W_2$ ordering reverses (flow $0.148$ vs.\ A-DCVAE $0.156$), so the accuracy advantage is scale-dependent, and we cite the canonical scale that matches our checkpoints. A patience-matched rerun rules out under-training as the explanation for the flow's canonical-scale gap: giving it five times more patience made it \emph{worse} (more over-dispersed), not better.

\subsection{Isolating the three guarantees}
\label{sec:ablation}

Finally, we isolate the three ingredients the theory assigns distinct roles (\Secref{subsec:theory}), ablating each noise source in turn over $5$ seeds at canonical scale (Figure~\ref{fig:ablation}). Removing \emph{condition} noise (N1) raises the empirical contraction constant from $\kappa = 0.41$ to $0.71$, toward the stability limit $\kappa = 1$---the predicted loss of the Jacobian/Tikhonov control; chain divergence stays zero (the state space is bounded), so the effect is on mixing \emph{rate}, not blow-up. Replacing the stochastic kernels with their deterministic mean maps (N4 off) collapses the along-fiber spread from $0.72$ to $0.000$ for the full model---regression-to-the-mean, reproduced on demand---confirming that injection noise is what averts collapse. Consistent with our own altitude assessment, \emph{target} noise (N2) shows no isolated benefit here (across-fiber $mDI$ error $\approx$ full), matching its status as a background principle rather than a load-bearing guarantee. The self-consistency certificate $\varepsilon_{\mathrm{inc}}$ stays ${\approx}0.03$ across variants, indicating $\pi^\ast$ is near self-consistent. One richer detail at canonical scale: when $\kappa$ is far from zero (N1 off), even the deterministic mean-map ping-pong no longer fully collapses (spread $0 \to 1.72$) within the step budget, reinforcing that the contraction knob (N1) and the non-collapse knob (N4) are genuinely separate.

\begin{figure}[t]
    \centering
    \includegraphics[width=\textwidth]{figures/figure_b7_ablation.pdf}
    \caption{\textbf{Three-guarantee isolation ablation} (mean $\pm$ std, $5$ seeds, canonical scale). Left to right, top then bottom: contraction $\kappa$ and chain divergence (N1 / stability), across-fiber $mDI$ error (N2 / direction), along-fiber spread stochastic vs.\ deterministic (N4 / non-collapse), self-consistency $\varepsilon_{\mathrm{inc}}$, and $\pi^\ast$ mass in the reference HPD-$95$. Removing condition noise raises $\kappa$; removing injection noise collapses the along-fiber spread to zero.}
    \label{fig:ablation}
\end{figure}



\begin{figure}[h]
\begin{center}
\fbox{\rule[-.5cm]{0cm}{3.5cm} \rule[-.5cm]{6cm}{0cm}}
\end{center}
\caption{\textbf{Posterior Predictive Check:} Forward ODE integration using samples drawn from the two peaks of our estimated bimodal posterior. Both parameter regimes reconstruct the sparse empirical observation track (dots) with exact precision, proving the model has decoded the true underlying system ambiguity.}
\label{fig:predictive_check}
\end{figure}

\section{Analysis: Expressiveness vs. Robustness Trade-off}
\label{sec:analysis}
To stress-test our model, we evaluate inference degradation by incrementally injecting Gaussian noise into the test observations $\mX_{\mathrm{obs}}$. As illustrated in Figure~\ref{fig:noise}, this experiment uncovers a compelling trade-off between a model's local expressiveness and its vulnerability to Out-of-Distribution (OOD) perturbation.

\begin{figure}[h]
\begin{center}
\fbox{\rule[-.5cm]{0cm}{3.5cm} \rule[-.5cm]{6cm}{0cm}}
\end{center}
\caption{\textbf{Noise Sensitivity Profile:} Deterministic baselines maintain a high-bias, flat-line error response. In contrast, generative architectures exhibit high-variance spiking curves once observation noise pushes inputs off the clean manifold.}
\label{fig:noise}
\end{figure}

The high sensitivity of the Dual CVAE originates directly from its superior performance: it maps a razor-thin, precise geometric manifold of the physics equations. When observation noise pushes inputs off-manifold, the generative alignment falters, leading to variance spikes. This reveals that structural precision introduces sensitivity to noise, defining a distinct bias-variance frontier.

\section{Conclusion \& Future Work}
\label{sec:conclusion}
We presented a probabilistic decoupled framework for nonlinear inverse problems that keeps the alternating operator of deterministic decoupled estimators but changes what it converges to, replacing global spectral contraction with local, noise-driven guarantees: existence of a unique invariant law by ergodicity and steering toward the true posterior by denoising. By sampling conditionals over the complete physical state space, it recovers structured multimodal posteriors along non-identifiable manifolds where mean-seeking estimators collapse.

\paragraph{Limitations.}
Our guarantees and experiments target non-identifiability that is \emph{structural}---a low-dimensional fiber induced by missing states. Strong \emph{practical} non-identifiability (near-flat likelihoods without a clean fiber), posteriors with three or more well-separated modes, and observation noise that pushes inputs off the training manifold all stress the method; the noise sensitivity documented in \Secref{sec:analysis} is a direct symptom of the last. The steering guarantee is only as good as the simulator: it bounds the gap to the posterior \emph{under the assumed forward model}, not to reality, so the Sim2Real gap remains open---this is why we frame domain adaptation as future work rather than as a contribution or a demonstrated capability. Finally, the coverage--sharpness trade-off of \Secref{sec:experiments} means our intervals are wide by construction along unidentifiable directions; this width is faithful rather than free, and aggregate interval metrics can undervalue it. Bridging the Sim2Real gap via noise-aware and domain-adaptive training is the natural next step.

\bibliography{iclr2026_conference}
\bibliographystyle{iclr2026_conference}

\appendix

\section{Proofs of Theoretical Properties}
\label{app:proofs}

We provide the proofs for the three properties stated in \Secref{subsec:theory}. Throughout, we
argue symmetrically over the two decoders of A-DCVAE---the Hidden-State CVAE decoder $g_\phi$ and
the Parameter CVAE decoder $g_\psi$---and drop the sub/superscripts when a statement holds for both.
We write the generic decoder as $g(\vu, \xi, \vz)$, where $\vu = \mX_{\mathrm{obs}}$ is the fixed
context, $\xi$ is the conditioning variable subject to noise, and $\vz$ is the CVAE latent. We assume
$g$ is $C^2$ in $\xi$ and that expectations over $\vz \sim \mathcal{N}(\vzero, \mI)$ are finite.

\subsection{Proof of Theorem~\ref{thm:contraction} (Condition noise $\Rightarrow$ local Jacobian regularization)}
\label{app:proof_thm1}

Let the condition be perturbed by $\check{\xi} = \xi + \boldsymbol{\epsilon}_c$ with
$\boldsymbol{\epsilon}_c \sim \mathcal{N}(\vzero, \sigma_c^2 \mI)$, while the decoder is trained to
reconstruct the clean target $Y$. A second-order Taylor expansion of the decoder about the clean
condition $\xi$ gives
\begin{equation}
    g(\vu, \xi + \boldsymbol{\epsilon}_c, \vz)
    = g(\vu, \xi, \vz) + \nabla_{\xi} g\, \boldsymbol{\epsilon}_c
      + \tfrac{1}{2}\, \boldsymbol{\epsilon}_c^\top \nabla_{\xi}^2 g\, \boldsymbol{\epsilon}_c
      + \mathcal{O}(\|\boldsymbol{\epsilon}_c\|^3).
\end{equation}
Substituting into the reconstruction term of the ELBO and taking the expectation over
$\boldsymbol{\epsilon}_c$, the linear term vanishes because $\E[\boldsymbol{\epsilon}_c] = \vzero$,
and the leading quadratic contribution is
\begin{equation}
    \E_{\boldsymbol{\epsilon}_c}\!\left[ \boldsymbol{\epsilon}_c^\top (\nabla_{\xi} g)^\top (\nabla_{\xi} g)\, \boldsymbol{\epsilon}_c \right]
    = \sigma_c^2\, \Tr\!\big( (\nabla_{\xi} g)^\top (\nabla_{\xi} g) \big)
    = \sigma_c^2\, \big\| \nabla_{\xi} g \big\|_F^2 .
\end{equation}
Taking the expectation over the latent prior and using that, for a sufficiently expressive network,
the mean residual $\E[Y - g]$ is small so the Hessian cross-term is higher order, the objective
reduces asymptotically to
\begin{equation}
\label{eqn:tikhonov}
    \gL^{\mathrm{noisy}} \;\approx\; \gL^{\mathrm{clean}}
    \;+\; \sigma_c^2\, \E_{\vz}\!\left[ \big\| \nabla_{\xi} g(\vu, \xi, \vz) \big\|_F^2 \right]
    \;+\; \mathcal{O}(\sigma_c^4).
\end{equation}
Equation~\eqref{eqn:tikhonov} is exactly a Tikhonov penalty on the expected condition-Jacobian. We
emphasize that this \emph{suppresses} $\|\nabla_{\xi} g\|_F$ rather than certifying a global
Lipschitz bound. Its dynamical consequence is dissipativity: for the composed transition, the
expected one-step change of $\| \nabla g \|_F$ is negative outside a bounded region, so the map acts
as a Lyapunov drift that keeps the chain within a bounded set instead of amplifying perturbations.
Combined with the compact boundary projections $\mathcal{B}_y, \mathcal{B}_p$, this prevents the
exponential error amplification of Issue~1. $\hfill\square$

\subsection{Proof of Theorem~\ref{thm:score} (Target noise $\Rightarrow$ denoising score matching)}
\label{app:proof_thm2}

Let the target be perturbed by $\tilde{Y} = Y + \boldsymbol{\epsilon}_t$ with
$\boldsymbol{\epsilon}_t \sim \mathcal{N}(\vzero, \sigma_t^2 \mI)$. By the calculus of variations, the
reconstructor minimizing the expected squared error is the posterior mean,
\begin{equation}
    g^*(\tilde{Y}) = \E\big[ Y \mid \tilde{Y}, \vu, \vtheta \big]
    = \int Y\, p\big(Y \mid \tilde{Y}, \vu, \vtheta\big)\, dY .
\end{equation}
By Bayes' rule, $p(Y \mid \tilde{Y}) \propto \mathcal{N}(\tilde{Y} \mid Y, \sigma_t^2 \mI)\,
p_{\mathrm{data}}(Y \mid \vu, \vtheta)$. Using the Gaussian identity
$\nabla_{\tilde{Y}} \mathcal{N}(\tilde{Y} \mid Y, \sigma_t^2 \mI)
= -\tfrac{\tilde{Y}-Y}{\sigma_t^2}\, \mathcal{N}(\tilde{Y} \mid Y, \sigma_t^2 \mI)$ to isolate
$Y\,\mathcal{N}$, substituting into the integral, and applying Leibniz's rule yields Tweedie's
formula \citep{efron2011tweedie},
\begin{equation}
    g^*(\tilde{Y}) - \tilde{Y}
    \;=\; \sigma_t^2\, \nabla_{\tilde{Y}} \log p_{\mathrm{data}}\big(\tilde{Y} \mid \vu, \vtheta\big).
\end{equation}
Hence the decoder's residual equals the score of the smoothed conditional density. Injecting noise
at inference and applying $g^*$ therefore realizes one step of a Langevin drift toward high-density
regions of $\mathcal{M}$, i.e.\ an attractor field that resolves Issue~2. $\hfill\square$

\subsection{Proof of Theorem~\ref{thm:ergodicity} (Ergodicity and score-aligned convergence)}
\label{app:proof_thm3}

Consider the alternating kernel $\mathcal{T}$ of \eqref{eqn:kernel}, obtained by composing the two
infinite Gaussian mixture kernels.

\textbf{(i) Unique invariant measure.} Each sub-kernel integrates an isotropic Gaussian against the
latent prior, so its transition density is strictly positive, $k(x, y) > 0$, on the compact domain
$\Omega$ enforced by $\mathcal{B}_y, \mathcal{B}_p$. Strict positivity implies the chain is
$\varphi$-irreducible, aperiodic, and Harris recurrent. By Theorem~\ref{thm:contraction}, the
Tikhonov-induced dissipativity supplies a Lyapunov drift condition of Meyn--Tweedie type \citep{meyn2009markov}. Together
these give geometric ergodicity: there exist $\rho \in (0,1)$ and $C < \infty$ such that
$W_2\big(\mathcal{T}^k(\mu_0, \cdot),\, \pi^*\big) \le C\, \rho^k$ in the $2$-Wasserstein metric, for
a unique invariant distribution $\pi^*$.

\textbf{(ii) Score-aligned target and no collapse.} Substituting the optimal reconstructor of
Theorem~\ref{thm:score} into the update shows that each step is an Euler--Maruyama discretization of
the unadjusted Langevin dynamics \citep{welling2011bayesian}
\begin{equation}
    \mX_{\mathrm{hid}}^{(k+1)} \leftarrow \mX_{\mathrm{hid}}^{(k)}
    + \alpha\, \nabla \log p_{\mathrm{data}}\big(\mX_{\mathrm{hid}}^{(k)}\big) + \sqrt{2\alpha}\, \boldsymbol{\epsilon},
    \qquad \alpha \propto \sigma_t^2 ,
\end{equation}
whose stationary law (via the Fokker--Planck equation) is the score-defining density. Because a fresh
latent $\vz$ is drawn at every step, $\mathcal{T}$ is an infinite mixture with the capacity to place
mass along an entire multi-modal (e.g.\ crescent-shaped) support rather than at a single mean.
Consequently $\pi^*$ aligns with the true non-identifiable manifold and the chain bypasses the
mode collapse of deterministic solvers. $\hfill\square$

\paragraph{Remark (scope of the guarantees).}
Theorem~\ref{thm:contraction} yields a Jacobian penalty and dissipativity, not a certified global
contraction; the invariant law $\pi^*$ is the stationary distribution of the learned stochastic
operator and approximates $p_{\mathrm{data}}$ up to the score-matching and $\epsilon$-incompatibility
errors of the trained decoders. These are the intended, and empirically supported, senses of the
three guarantees.

\end{document}